\BourbakiUnnumberedChapter{参考文献}{参考文献}

\begin{enumerate}
\def\labelenumi{\arabic{enumi}.}
\item
  O. 诺伊格鲍尔，《古代数学史讲义》，第I卷：前希腊数学，柏林（施普林格），1934年。
\item
  《欧几里得原本》，5卷，J. L. 海伯格编，莱比锡（托伊布纳），1883-88年。
\end{enumerate}

2 bis. T. L. 希思，《欧几里得〈原本〉十三卷》，\ldots，3卷，剑桥，1908年。

\begin{enumerate}
\def\labelenumi{\arabic{enumi}.}
\setcounter{enumi}{2}
\item
  H. 卡尔达诺，《著作集》，里昂，1663年。
\item
  R. 邦贝利，《代数学》，博洛尼亚（G. 罗西），1572年。
\item
  F. 维埃特，《数学著作集》，\ldots，莱顿（埃尔泽菲尔），1646年。
\item
  A. 吉拉尔，《代数学中的新发现》，阿姆斯特丹，1629年。
\item
  R. 笛卡尔，《几何学》，Fr. 范·斯霍滕拉丁文译本，第2版，2卷，阿姆斯特丹（埃尔泽菲尔），1659-61年。
\item
  G. W. 莱布尼茨，《数学著作》，C. I. 格哈特编，7卷，柏林-哈雷（阿舍尔-施密特），1849-63年。
\item
  《戈特弗里德·威廉·莱布尼茨与数学家通信集》，C. I. 格哈特编，第I卷，柏林（迈尔与米勒），1899年。
\item
  L. 欧拉，《全集》（1），第VI卷，柏林-莱比锡（托伊布纳），1921年：a)《论方程的根的形式》\ldots，第1-19页；b)《论任意次方程的求解》，第170-196页。
\item
  J.-L. 拉格朗日，《著作集》，第III卷，巴黎（戈蒂埃-维拉尔），1869年：a)《关于方程的代数求解的思考》，第205-421页；b)《论方程虚根的形式》，第479页。
\item
  A. 范德蒙德，《关于方程求解的论文》，皇家科学院历史，1771年，巴黎（1774年），第365-416页。
\item
  C. F. 高斯，《全集》，第I-X卷，哥廷根，1863-1923年。
\item
  N. H. 阿贝尔，《著作集》，2卷，西洛和李编，克里斯蒂安尼亚，1881年。
\item
  E. 伽罗瓦，《数学著作集》，巴黎（戈蒂埃-维拉尔出版社），1897年。
\item
  A.-L. 柯西，《全集》（第一辑），第X卷，巴黎（戈蒂埃-维拉尔出版社），1897年，第312页及第351页。
\item
  J. 刘维尔，《关于值既非代数数、甚至不可化为代数无理数的极广泛类别的量》，数学杂志（第一辑），第XVI卷（1851年），第133页。
\item
  R. 戴德金，《数学全集》，3卷，不伦瑞克（菲韦格出版社），1932年。
\item
  L. 克罗内克：a)《代数量算术理论纲要》，克雷尔杂志，第XCII卷（1882年），第1-122页（=《全集》，第II卷，莱比锡（托伊布纳出版社），1897年，第245-387页）；b)《一般算术的一个基本定理》，克雷尔杂志，第C卷（1887年），第490-510页（=《全集》，第III卷 \(_{1}\) ，莱比锡（托伊布纳出版社），1899年，第211-240页）。
\item
  H. 韦伯，《关于伽罗瓦方程理论的一般基础的研究》，数学年鉴，第XLIII卷（1893年），第521-544页。
\item
  G. 韦罗内塞，《几何学基础》，帕多瓦，1891年。
\item
  K. 亨泽尔，《代数数论》，莱比锡-柏林（托伊布纳出版社），1908年。
\item
  E. 施泰尼茨，《域的代数理论》，克雷尔杂志，第CXXXVII卷（1910年），第167-309页。
\item
  W. 克鲁尔，《无限代数扩张的伽罗瓦理论》，数学年鉴，第C卷（1928年），第687页。
\item
  E. 阿廷，《伽罗瓦理论》\ldots，安阿伯，1946年。
\item
  A. 韦伊，《代数几何基础》，美国数学学会论文集丛刊，第XXIX卷，纽约，1946年。
\end{enumerate}

